% Problem record op_1576e5f52f7ff3b2. This TeX file is derived from
% database/problems_json/op_1576e5f52f7ff3b2.json by scripts/sync-tex.mjs; edit the JSON record
% and rerun the script rather than editing this file.
\section{Is bipartite Quantum Max-Cut in BPP?}
\label{sec:op_1576e5f52f7ff3b2}

\paragraph{Problem.}
Is the following bipartite Quantum Max-Cut promise problem in
$\mathrm{BPP}$? Given a bipartite graph $G=(V,E)$ with $|V|=n$,
polynomially bounded nonnegative rational weights $w_{ij}$, and rational
thresholds $a<b$ with $b-a\geq1/\operatorname{poly}(n)$, define
\begin{equation}
  H_G:=\sum_{\{i,j\}\in E}w_{ij}(X_iX_j+Y_iY_j+Z_iZ_j),
  \qquad E_0:=\lambda_{\min}(H_G).
  \label{eq:bqmc-ground-energy}
\end{equation}
Here $X_i,Y_i,Z_i$ are Pauli operators on qubit $i$.
For the energy in Eq.~\eqref{eq:bqmc-ground-energy}, distinguish
$E_0\leq a$ from $E_0\geq b$, promised one holds, using a randomized
classical algorithm polynomial in the input length and correct with
probability at least $2/3$.

\subsection*{Status}

Unsolved

\subsection*{Source}

The remaining classical part of Open question 3.4 in Gharibian's
\emph{The 7 faces of quantum NP}, Section 3, arXiv page 7
\sourcecite{ref:bqmc-gharibian}{Gha24}.
The original question asks for the complexity of bipartite QMC;
the present formulation focuses on $\mathrm{BPP}$ membership following
the 2026 $\mathrm{BQP}$ upper bound \sourcecite{ref:bqmc-rt}{RT26}.

\subsection*{Progress}

\begin{itemize}
  \item Earlier upper bound: conjugating one bipartition by Pauli $Y$ makes
$H_G$ stoquastic, placing bipartite QMC in $\mathrm{StoqMA}$
(summarized in Section 3) \sourcecite{ref:bqmc-gharibian}{Gha24}.

  \item July 2026: Rayudu--Takahashi prove a Lee-Yang spectral gap at least
$h/4$ under field strength $h>0$, enabling adiabatic ground-energy
estimation to inverse-polynomial additive error and establishing
$\mathrm{BQP}$ membership (Theorem 11; Section 1.1)
\sourcecite{ref:bqmc-rt}{RT26}.

  \item Independently, Bravyi--Gosset--Liu--Wong give a
$\operatorname{poly}(n,J,1/\epsilon)$-time quantum algorithm for
additive-$\epsilon$ ground-energy estimation on weighted bipartite graphs,
where $J=\max_{\{i,j\}\in E}w_{ij}$ (Corollary 2)
\sourcecite{ref:bqmc-bglw}{BGLW26}.
\end{itemize}

\subsection*{References}

\begin{enumerate}[label={},leftmargin=4.8em,labelsep=0.5em]
  \footnotesize
  \item[\textup{[Gha24]}]\label{ref:bqmc-gharibian}
  S. Gharibian, "Guest Column: The 7 faces of quantum NP,"
\emph{ACM SIGACT News} \textbf{54}(4), 54--91 (2024).
\href{https://doi.org/10.1145/3639528.3639535}{doi:10.1145/3639528.3639535};
\href{https://arxiv.org/abs/2310.18010}{arXiv:2310.18010} (2023 preprint).

  \item[\textup{[RT26]}]\label{ref:bqmc-rt}
  C. Rayudu and J. Takahashi, "Spectral gap of Lee-Yang Hamiltonians,"
arXiv preprint (July 2026).
\href{https://doi.org/10.48550/arXiv.2607.10765}{doi:10.48550/arXiv.2607.10765};
\href{https://arxiv.org/abs/2607.10765v1}{arXiv:2607.10765v1}.

  \item[\textup{[BGLW26]}]\label{ref:bqmc-bglw}
  S. Bravyi, D. Gosset, Y. Liu, and B. Wong,
"Efficient quantum algorithm for Heisenberg spin systems,"
arXiv preprint (July 2026).
\href{https://doi.org/10.48550/arXiv.2607.14401}{doi:10.48550/arXiv.2607.14401};
\href{https://arxiv.org/abs/2607.14401v1}{arXiv:2607.14401v1}.
\end{enumerate}

\subsection*{Comment}

Membership in $\mathrm{BPP}$ remains open: the quantum upper bound
leaves classical complexity unresolved \sourcecite{ref:bqmc-rt}{RT26}.
Here $\mathrm{BPP}$ and $\mathrm{BQP}$ denote their promise-problem
versions. The target is inverse-polynomial additive precision.

\subsection*{Field}

Quantum algorithm

\subsection*{Topic}

Computational complexity and computability; Hamiltonian complexity; Quantum Max-Cut

\subsection*{ID}

\texttt{op\_1576e5f52f7ff3b2}
